% Add "answers" option to print the answers
\documentclass[addpoints]{exam}

\usepackage[utf8]{inputenc}
\usepackage{calc}
\usepackage[dvipsnames]{xcolor}
\usepackage{ifthen}
\usepackage{tcolorbox}

% \makebox[\textwidth]{Name:\enspace\hrulefill}
%
% \vspace{5mm}
%\makebox[\textwidth]{Instructor’s name:\enspace\hrulefill}

\definecolor{MyLightGray}{gray}{0.60}

\newcommand{\truefalse}
{		\begin{checkboxes}
			 \choice True
			 \choice False. Explanation:
		\end{checkboxes}
		\vspace{10pt}
}


% Check if 'answers' is an option
\newcommand{\ifanswers}[2]{%
  \ifthenelse{\boolean{exam@answers}}{#1}{#2}%
}



\newcommand{\answerspace}[1]
{
	\ifprintanswers
	\else
		\setlength{\mylength}{10pt}%
		% \vspace{#1\mylength}
		% Make the grey box
		\par\nobreak\noindent
		\makebox[\linewidth][l]{%
			% set inner padding to zero, so only the frame shows
			{\setlength{\fboxsep}{0pt}%
			% draw a color frame, white (or transparent) inside
			\fcolorbox{MyLightGray}{white}{%
			% the inner box has zero width content but fixed height
			\parbox[t][#1\mylength][c]{%
				% subtract twice the frame‐thickness from the full width
				\dimexpr\linewidth-2\fboxrule\relax
			}{\strut}%
			}%
			}%
		}%
		\par
		\vspace{3\mylength}

  	\fi
}
\newlength{\mylength}

\newcommand{\code}[1]{{\texttt{#1}}}


%\usepackage{listings}
\usepackage{color}
\usepackage{lastpage}
\usepackage{verbatim}
\usepackage{fancyvrb}
\usepackage{graphicx}
\usepackage{float}
\usepackage{minted}
\usepackage{multicol}

\newcommand{\includecode}[2][c]{\lstinputlisting[caption=, escapechar=, language=#1]{#2}}
\RecustomVerbatimCommand{\VerbatimInput}{VerbatimInput}{}

\renewcommand{\solutiontitle}{\noindent\textbf{Solution sketch: }}

\definecolor{SolutionColor}{rgb}{0.8,0.9,1}

\cfoot{Page \thepage\ of \pageref{LastPage}}

\rhead{Initials: \rule{20mm}{0.4pt}}



%%%%%%%%%%%%%%%%%%%%%%%%%%%%%%%%%%%%%%%%
\begin{document}
\fvset{numbers=left,numbersep=3pt}
\thispagestyle{empty}

\begin{center}
	{\LARGE CS 408 Fall 2026 -- Data-Flow Coverage Quiz 2.0 \\}
\fbox{\fbox{\parbox{5.5in}{\centering
	\begin{itemize}
        \item Try to come up with a solution on your own first, but
          after that talk to your friends and neighbors.

        \item Full points will be rewarded for \textbf{any} (reasonable)
          attempt.
        \end{itemize}
	}}}
\end{center}

\vspace{5mm}
\makebox[\textwidth]{\textbf{Name (please print)}: \rule{100mm}{0.8pt} }

\vspace{10mm}
\makebox[\textwidth]{\textbf{PUID}: \rule{100mm}{0.8pt} \hspace{35mm}}


\begin{questions}

\begin{minted}{java}
static int compare(int a,
                   int b) {
  int x = 0;   // L3
  if (a > 0) { // C4
    x = 1;     // L5
  }
  if (b > 0) { // C7
    x = 2;     // L8
  }
 return x;     // L10
}
\end{minted}

\question Write down a modified version of the
  \mintinline{java}{compare} function shown above, whre satisfying
  100\% all-uses coverage requires strictly more tests than satisfying
  100\% all-defs coverage. Try to keep your edits as small as
  possible.

    \answerspace{26} %
    \begin{solution} %
      $def(\texttt{L2}) = \{y\}$, \;
      $def(\texttt{L4}) = \{y\}$, \;
      $def(\texttt{L6}) = \{y\}$, \;
      $def(\texttt{C3}) = \emptyset$, \;
      $def(\texttt{L7}) = \emptyset$.

      $use(\texttt{L7}) = \{y\}$ (a c-use); every other node has an
      empty use set.

      $use((\texttt{C3},\texttt{BB4})) =
       use((\texttt{C3},\texttt{BB6})) = \{x\}$ --- the p-use of
      \mintinline{python}{x} sits on \emph{both} edges out of the
      decision, not on \texttt{C3} itself.
    \end{solution}

\end{questions}

\end{document}
